% AUTO-GENERATED by tools/usc_extract.py — do not edit by hand
\begingroup\small
\begin{historicalsource}[{\textbf{Primary text} --- Title~18, \S~926B (OLRC via \texttt{nickvido/us-code}, tag \texttt{annual/2025})}]
\textit{Source:} \url{https://uscode.house.gov/view.xhtml?req=granuleid:USC-prelim-title18&num=0&edition=prelim}\par\medskip\textbf{§ 926B. Carrying of concealed firearms by qualified law enforcement officers}\\
\textbf{(a)} Notwithstanding any other provision of the law of any State or any political subdivision thereof, an individual who is a qualified law enforcement officer and who is carrying the identification required by subsection (d) may carry a concealed firearm that has been shipped or transported in interstate or foreign commerce, subject to subsection (b).

\textbf{(b)} This section shall not be construed to supersede or limit the laws of any State that—

(1) permit private persons or entities to prohibit or restrict the possession of concealed firearms on their property; or

(2) prohibit or restrict the possession of firearms on any State or local government property, installation, building, base, or park.

\textbf{(c)} As used in this section, the term “qualified law enforcement officer” means an employee of a governmental agency who—

(1) is authorized by law to engage in or supervise the prevention, detection, investigation, or prosecution of, or the incarceration of any person for, any violation of law, and has statutory powers of arrest or apprehension under \href{https://uscode.house.gov/view.xhtml?req=granuleid:USC-prelim-title10-section807/b&num=0&edition=prelim}{section 807(b) of title 10}, United States Code (article 7(b) of the Uniform Code of Military Justice);

(2) is authorized by the agency to carry a firearm;

(3) is not the subject of any disciplinary action by the agency which could result in suspension or loss of police powers;

(4) meets standards, if any, established by the agency which require the employee to regularly qualify in the use of a firearm;

(5) is not under the influence of alcohol or another intoxicating or hallucinatory drug or substance; and

(6) is not prohibited by Federal law from receiving a firearm.

\textbf{(d)} The identification required by this subsection is the photographic identification issued by the governmental agency for which the individual is employed that identifies the employee as a police officer or law enforcement officer of the agency.

\textbf{(e)} As used in this section, the term “firearm”—

(1) except as provided in this subsection, has the same meaning as in section 921 of this title;

(2) includes ammunition not expressly prohibited by Federal law or subject to the provisions of the National Firearms Act; and

(3) does not include—

  (A) any machinegun (as defined in section 5845 of the National Firearms Act);

  (B) any firearm silencer (as defined in section 921 of this title); and

  (C) any destructive device (as defined in section 921 of this title).

\textbf{(f)} For the purposes of this section, a law enforcement officer of the Amtrak Police Department, a law enforcement officer of the Federal Reserve, or a law enforcement or police officer of the executive branch of the Federal Government qualifies as an employee of a governmental agency who is authorized by law to engage in or supervise the prevention, detection, investigation, or prosecution of, or the incarceration of any person for, any violation of law, and has statutory powers of arrest or apprehension under \href{https://uscode.house.gov/view.xhtml?req=granuleid:USC-prelim-title10-section807/b&num=0&edition=prelim}{section 807(b) of title 10}, United States Code (article 7(b) of the Uniform Code of Military Justice).
\end{historicalsource}
\endgroup
